\section{Global convergence: homogeneous and mixed geometry}\label{art:global}

Toward a minimizer, a homogeneous term has an exact first-order model;
a smooth convex term adds only a quadratic remainder. We turn this inward
accuracy into the secant bound needed by the comparison energy, then apply
it to centered sums and coercive homogeneous objectives in any finite
dimension.

\subsection{Why inward accuracy controls a nonradial algorithm}

A gradient jump can make the positive correction in
\eqref{art:energy-identity} large. We control its pairing with the direction
to a minimizer by asking that first-order prediction remain quadratically
accurate on short inward segments. This asks nothing about transverse
smoothness, and it does not require BFGS to take radial steps. The comparison
segments are an analytical device; the actual steps remain those in
\eqref{art:update}.

Fix a minimizer $x_*$, write $F_*=F(x_*)$, and let
$K=\{x:F(x)\le b\}$ be compact, with $b>F_*$. For differentiable $x,z\in K$,
set
\[
 \begin{gathered}
 r_x=x-x_*,\qquad g=\nabla F(x),\qquad h=\nabla F(z),\\
 T=(h-g)^T(z-x),\qquad N=(h-g)^Tr_z.
 \end{gathered}
\]
Here $T$ is the quantity $t$ of Section~\ref{art:energy}, and $N$ is the
\emph{radial numerator}.
The radial secant condition is the following requirement on positive bands:
\begin{equation}\label{art:radial-condition}
 \begin{gathered}
 \text{for every }\eta>0\text{ there is }C_\eta\ge0\text{ such that}\\
 N^2\le C_\eta T
 \quad\text{whenever }F_*+\eta\le F(z)\le F(x),\quad T>0.
 \end{gathered}
\end{equation}
Here all such differentiable endpoint pairs are included, whether or not they
are eligible for the fixed line search. The inward model condition is
\begin{equation}\label{art:inward-condition}
 \begin{gathered}
 \text{for every }\eta>0\text{ there are }M_\eta\ge0,
       \ \theta_\eta\in(0,1/2]\text{ such that}\\
 0\le F(x-\tau r_x)-F(x)+\tau\nabla F(x)^Tr_x
       \le\frac{M_\eta}{2}\tau^2,\quad 0\le\tau\le\theta_\eta,
 \end{gathered}
\end{equation}
at every differentiable $x\in K$ with $F(x)\ge F_*+\eta$.
The inward point need not be differentiable. It remains in $K$ by convexity.
The parameter $\tau$ is a fraction of the distance to $x_*$, rather than a
Euclidean step length. Constants may depend on the positive value band:
no common bound is required as the band approaches the minimum. Thus
nonsmoothness at the solution is permitted, including the absolute-value
term in the benchmark.

\begin{theorem}[Inward models control radial secants]\label{art:radial-equivalence}
For finite convex $F$ and the compact sublevel $K$ above,
\eqref{art:inward-condition} implies \eqref{art:radial-condition}.
\end{theorem}
\begin{proof}
Choose $\Lambda,R>0$ bounding subgradient norms on $K$ and $\norm{x-x_*}$ there.
The subgradient bound was shown in the proof of Theorem~\ref{art:certificate}.
Use the supporting-plane error
\[
 D_F(z,x)=F(z)-F(x)-\nabla F(x)^T(z-x)\ge0.
\]
This is the usual Bregman expression \cite{BB}; only its supporting endpoint
needs to be differentiable. At differentiable endpoints,
$T=D_F(z,x)+D_F(x,z)$.

Suppose that \eqref{art:inward-condition} holds, and fix a positive band,
with constants $M,\theta$. If $N\ge0$, compare the supporting plane at $x$
with the inward model at $z$, both evaluated at $z-\tau r_z$, to obtain
\[
 \tau N\le D_F(z,x)+\tfrac12M\tau^2.
\]
If $N<0$, reverse the endpoints and use the inward model at $x$. The reversed
pair has radial numerator $(g-h)^Tr_x=T-N\ge|N|$, and
\[
 \tau(T-N)\le D_F(x,z)+\tfrac12M\tau^2.
\]
Thus there are $A\in[|N|,2\Lambda R]$ and
$D\in\{D_F(z,x),D_F(x,z)\}$ such that
$\tau A\le D+M\tau^2/2$ for $0\le\tau\le\theta$.
For $M>0$ and $A\le M\theta$, choose $\tau=A/M$, giving $A^2\le2MD$.
If $A>M\theta$, take $\tau=\theta$ to obtain
$D\ge A\theta/2$ and hence $A^2\le(4\Lambda R/\theta)D$.
When $M=0$, the choice $\tau=\theta$ gives an even stronger estimate.
Therefore
\begin{equation}\label{art:inward-to-radial}
 N^2\le \max\{2M,4\Lambda R/\theta\}\,T.\qquad\qed
\end{equation}
\end{proof}

The converse implication, including inward samples that are not
differentiable, is proved in Chapter~4 of the complete companion. The article
uses the sufficient direction just proved: objective modeling accuracy
controls every descending secant pair on the same positive band.

\begin{corollary}[Radial convergence]\label{art:radial-convergence}
Under either condition, every admissible run starting in $K$
converges in value to $F_*$ and in distance to $\arg\min F$. With a unique
minimizer it converges to $x_*$.
\end{corollary}
\begin{proof}
Take $W_\eta(x)=r_xr_x^T$. Convexity gives
$g^Tr_x\ge F(x)-F_*\ge\eta$, proving visibility with $m_\eta=\eta^2$.
By Theorem~\ref{art:radial-equivalence} we may assume \eqref{art:radial-condition};
it gives secant control on eligible pairs in $K_\eta$, since these have
$F(z)<F(x)$ and $T=t>0$.
Finally $r_z=r_x+s$ implies $\Pi_sr_z=\Pi_sr_x$, so transport holds with
$L_\eta=0$. Apply Theorem~\ref{art:certificate}.\qed
\end{proof}

Condition~\eqref{art:inward-condition} controls only inward first-order modeling error. In a smooth
region its infinitesimal expression is $r_x^T\nabla^2F(x)r_x$, rather than the
largest Hessian eigenvalue. A bound almost everywhere on that expression alone
would miss a jump of the derivative along a ray; the function-value condition
\eqref{art:inward-condition} includes such jumps. The condition is sufficient;
Theorem~\ref{art:nonradial-example} will exhibit convergence where
\eqref{art:radial-condition}, and hence by Theorem~\ref{art:radial-equivalence} also
\eqref{art:inward-condition}, fails on positive bands arbitrarily near the minimizer.

\subsection{A family theorem and the mixed benchmark}

A finite convex positively homogeneous function is sublinear: its value on
$(1-\tau)r$ is exactly $(1-\tau)$ times its value on $r$. Its angular
nonsmoothness consequently disappears from the inward remainder. The next
theorem permits an arbitrary locally $C^{1,1}$ convex perturbation, provided
the homogeneous part is centered at a minimizer of the sum. Centering is a
geometric hypothesis on the objective, not knowledge required by the
algorithm. Moreover, the two summands need not separately be minimized
at that center.

\begin{theorem}[Centered smooth-plus-sublinear objectives]\label{art:centered}
Let $F:\R^n\to\R$ be finite and convex, and suppose $x_*$ minimizes $F$.
Let $U$ be an open convex set, let $\psi$ be convex and locally $C^{1,1}$ on
$U$, and let $\gamma:\R^n\to\R$ be finite, convex and positively homogeneous
of degree one. Suppose
\[
 F(x)=\psi(x)+\gamma(x-x_*)\qquad(x\in U).
\]
Every admissible run whose initial sublevel is compact and contained in $U$
satisfies $F(x_k)\to F(x_*)$. The distance to $\arg\min F$ also tends to zero. If $x_*$ is the unique
minimizer, then $x_k\to x_*$. Neither original summand is required to be
minimized at $x_*$.
\end{theorem}
\begin{proof}
Let $K=\{F\le F(x_0)\}$. Since $g_0\ne0$, convexity implies
$F(x_0)>F(x_*)$. On compact convex $K\subset U$, local $C^{1,1}$ regularity
provides a uniform Lipschitz constant $L$ for $\nabla\psi$: take a finite
cover by neighborhoods carrying local constants and subdivide any segment in
$K$ into pieces shorter than a Lebesgue number of that cover. Summing the
local bounds over the pieces gives the largest of those constants as a
Lipschitz constant along the segment. Integrating the gradient along an
inward segment consequently gives
\[
 \left|\psi(x-\tau r_x)-\psi(x)
             +\tau\nabla\psi(x)^Tr_x\right|
 \le\tfrac12L\tau^2\norm{r_x}^2.
\]
At a differentiable point of $F$ in $U$, the representation and smoothness of
$\psi$ imply differentiability of $\gamma$ at $r_x$. Differentiating its
homogeneity identity gives $\nabla\gamma(r_x)^Tr_x=\gamma(r_x)$.
Its contribution to the inward remainder is therefore exactly zero.
The upper bound in \eqref{art:inward-condition} follows with
$M=L\max_{x\in K}\norm{r_x}^2$ and $\theta=1/2$ for every band; its lower bound
follows from convexity of $F$.
Corollary~\ref{art:radial-convergence} proves the claim.\qed
\end{proof}


The compact initial sublevel supplies uniform constants and traps the accepted
iterates. When the representation is global and all sublevels are compact,
the theorem applies from every differentiable nonstationary starting point,
for every initial SPD matrix and every fixed valid Armijo--weak-Wolfe pair.
Neither a lower Hessian bound for the smooth term nor differentiability of
the full objective away from the minimizer is required. Without uniqueness,
the stated conclusion is value and solution-set distance convergence;
convergence to one particular minimizer is not asserted.

\begin{corollary}[The quadratic--absolute-value benchmark]\label{art:mixed-benchmark}
For $F(u,v)=u^2+|v|$, every infinite admissible full-BFGS run, with arbitrary
$H_0\succ0$ and any fixed $0<c_1<c_2<1$, satisfies
\[
 F(u_k,v_k)\longrightarrow0,\qquad (u_k,v_k)\longrightarrow(0,0).
\]
\end{corollary}
\begin{proof}
Take $\psi(u,v)=u^2$, $\gamma(u,v)=|v|$, $x_*=(0,0)$ and $U=\R^2$.
The objective is finite and convex, its sublevels are closed and bounded,
and its minimizer is unique. Equation~\eqref{art:mixed-inward-remainder}
shows the exact inward error. Theorem~\ref{art:centered} applies.\qed
\end{proof}

This strengthens Guo--Lewis's subsequence approach to the line $v=0$
\cite[Corollary~4.4]{GL18} to the whole-sequence convergence to the origin
they left open, under the admissible full-BFGS convention of this paper.

\subsection{Coercive homogeneous objectives, linear maxima, and gauges}

Pure homogeneous geometry is the limiting case with no smooth remainder.
Coercivity here has a clear purpose: it prevents escape along a direction of
zero or negative homogeneous value and converts convergence in value into
convergence to the origin. It imposes no smoothness or finite-piece condition.

\begin{theorem}[Coercive homogeneous objectives]\label{art:homogeneous}
Let $F:\R^n\to\R$, $n\ge1$, be finite, convex, positively homogeneous of
degree one, and coercive. Then zero is its unique minimizer and every
infinite admissible full-BFGS run satisfies
\[
 F(x_k)\longrightarrow0,\qquad x_k\longrightarrow0.
\]
The initial point, initial SPD matrix and fixed valid line-search constants
are arbitrary within admissibility.
\end{theorem}
\begin{proof}
Homogeneity gives $F(0)=0$. If $F(x)\le0$ for some $x\ne0$, then
$F(tx)=tF(x)\le0$ as $t\to\infty$, contradicting coercivity. Thus
$F(x)>0$ off zero. Continuity on the unit sphere gives
\[
 \mu=\min_{\norm{x}=1}F(x)>0,
 \qquad F(x)\ge\mu\norm{x}.
\]
All sublevels are compact. Theorem~\ref{art:centered}, with $x_*=0$, $U=\R^n$, $\psi=0$ and
$\gamma=F$, gives convergence in value, and the displayed lower bound gives
convergence of the iterates. Alternatively, the inward remainder is identically
zero, so the radial convergence corollary applies directly.\qed
\end{proof}

\begin{corollary}[Compact homogeneous models]\label{art:homogeneous-examples}
The conclusion of Theorem~\ref{art:homogeneous} holds for every norm, every
finite convex degree-one gauge with bounded unit sublevel, and every finite
linear maximum
\[
 F(x)=\max_{1\le i\le m}a_i^Tx
\]
whose zero sublevel is bounded. These are whole-iterate conclusions from every
admissible start. The linear maximum has no affine intercepts.
\end{corollary}
\begin{proof}
For a norm, continuity and positivity give a positive minimum on the Euclidean
unit sphere. For either of the other models, boundedness of the stated
sublevel rules out a nonzero $x$ with $F(x)\le0$, since its full positive ray
would lie in that sublevel. The same sphere argument gives
$F(x)\ge\mu\norm{x}$, hence coercivity. Finite maxima of linear forms are
finite, convex and degree-one homogeneous, so the theorem applies.\qed
\end{proof}

The norm cases include $\norm{x}_1=\sum_i|x_i|$ and
$\norm{x}_\infty=\max_i|x_i|$. In dimension at least two their nonsmoothness
is nonisolated: the first has gradient jumps across coordinate hyperplanes;
the second has them where two nonzero coordinate magnitudes tie for the
maximum. For example, along $(t,t,0,\ldots,0)$ with $t>0$, the two active
coordinate forms of $\norm{\cdot}_\infty$ have distinct gradients. These
are concrete specializations of the norm corollary, with no separate
priority claim. The familiar
Euclidean norm is an already established consequence of Guo--Lewis's
smooth-extension analysis \cite[Corollary~4.2]{GL18}; the scalar absolute-value
model was analyzed by Lewis--Overton under their stated search conventions
\cite[Section~5.1]{LO}. No symmetry is needed for the gauge or max-linear
cases. Coercivity supplies this compact, whole-iterate route. The complete companion,
Chapter~7, separately proves the noncompact homogeneous value dichotomy
$F(x_k)\to0$ or $F(x_k)\to-\infty$, together with its bounded-below and
additional iterate consequences.
